\chapter{样式}

\section{转义字符}
注意“\#”是 \TeX 元字符，因此 C\# 要写作 \fn{C\bs\#}。
类似的“\textasciitilde”也是元字符，
要写作\mn{textasciitilde}。
\index{命令!textasciitilde@\mn{textasciitilde}}
这两个字符在URL中也经常出现，要特别小心。

C/C++代码的标识符中经常出现下划线“_”，例如 \fn{boost::shared_ptr}、
\fn{random_\linebreak[0]shuffle} 等。
为了避免每次都转义，可以使用 \fn{underscore} 宏包，
将下划线变为普通字符。
\index{宏包!underscore@\fn{underscore}}
注意这里 \fn{random_shuffle} 在行尾断字，
那么下划线之后不应该出现连字号“-”，
因此应写为 \verb|random_\linebreak[0]shuffle|。

\section{斜体}
按照排版规范，数学变量名应该用斜体，常量和单位应该用正体。
例如：$n$-body 问题、$E = m \mathrm{c}^2 $、5\textmu s。
快速排序 $n$ 个元素的数度的平均时间复杂度是 $O(n \log n)$。
通过 TCP 发送 $n$ 字节的消息，
接收方收到这 $n$ 个字节的事件可能性有 $2^{n-1}$ 种
\footnote{在 $n-1$ 个字节间隙中，依次插入 $0, 1, \ldots, n-1$ 个“隔板”，
一共有 $C_{n-1}^0 + C_{n-1}^1 + \cdots + C_{n-1}^{n-1} = 2 ^ {n-1} $ 种情况。}。

\section{列表}
\fn{enumitem}
\index{宏包!enumitem@\fn{enumitem}}

\section{章节标题}
章节标题无标点，因此 \S \ref{sec:whyTypesetting} 是错的。

\subsection{编号}
单个小节不编号，因此 \S \ref{subsec:noChineseItalic} 和本小节是错的。

\section{图表编号}

\section{脚注} %%%%%%%%%%%%%%%%%%%%%%%%%%%%%%

\subsection{编号}
\LaTeX默认是按章重置脚注编号，
这么如果重印的时候需要增加一个脚注，
势必会影响后续页码，这是 \mybooktitle 排版的一个教训。
因此我意识到脚注应该是按页重置，但是 \verb|\@addtoreset{footnote}{page}|
无效，必须使用 \fn{footmisc} 宏包的 \fn{perpage} 选项，
例如 \sfn{\bs usepackage[perpage]\{footmisc\}}。
\index{宏包!footmisc@\fn{footmisc}}

\LaTeX默认的脚注编号是数字，这有时会造成误解。
例如给长度单位pt添加脚注，正文中可能会出现“pt\textsuperscript{2}”，
让人误以为是面积单位，因此可以改用带圈数字。

\subsection{置底}
\centerline{\fbox{\includegraphics[page=1]{footnote-middle.pdf}}%
\quad\fbox{\includegraphics[page=2]{footnote-middle.pdf}}}

使用 \fn{footmisc} 宏包的 \fn{bottom} 选项之后，
脚注置底。\nopagebreak

\centerline{\fbox{\includegraphics[page=1]{footnote-bottom.pdf}}%
\quad\fbox{\includegraphics[page=2]{footnote-bottom.pdf}}}

\section{参考文献}
技术书籍不是学术著作，可不必使用 BibTeX 工具，直接排版参考文献即可。
